\documentclass[11pt,a4paper]{article}

\newcommand{\doctitle}{Event Frequency Model (MKM-EF-001)}
\newcommand{\docsubtitle}{Wolfram Oracle Benchmark of Annualisation, Tenor Probability, Compound Loss and Count Families}
\newcommand{\docversion}{0.2}
\newcommand{\docdate}{27-October-2026}
\newcommand{\docauthor}{David K Kelly}

% MKM Research Labs — Shared Document Header
% Usage: % MKM Research Labs — Shared Document Header
% Usage: % MKM Research Labs — Shared Document Header
% Usage: \input{../shared/mkm_header}  (before \begin{document})
%
% Required commands before \input:
%   \newcommand{\doctitle}{...}
%   \newcommand{\docsubtitle}{...}
%   \newcommand{\docversion}{1.0}
%   \newcommand{\docdate}{13-February-2026}
%   \newcommand{\docauthor}{David K Kelly}

% ─── Packages ───────────────────────────────────────────────────────────────
\usepackage[utf8]{inputenc}
\usepackage{amsmath,amsfonts,amssymb}
\usepackage{listings}
\usepackage{color,xcolor}
\usepackage{graphicx}
\usepackage{hyperref}
\usepackage[toc,page]{appendix}
\usepackage{tabto}
\usepackage{float}
\usepackage{comment}
\usepackage[en-GB]{datetime2}
\usepackage{setspace}
\usepackage{subcaption}
\usepackage{booktabs}
\usepackage{longtable}
\usepackage{adjustbox}
\usepackage{enumitem}
\usepackage{fancyhdr}
\usepackage{lastpage}

% ─── Page geometry ──────────────────────────────────────────────────────────
\onehalfspacing
\usepackage{geometry}
\geometry{a4paper, total={170mm,247mm}, left=20mm, top=25mm, bottom=25mm}

% ─── Headers and footers ───────────────────────────────────────────────────
\pagestyle{fancy}
\fancyhf{}
\fancyhead[L]{\small\textcolor{gray}{MKM Research Labs}}
\fancyhead[R]{\small\textcolor{gray}{\doctitle{} --- v\docversion}}
\fancyfoot[C]{\small\textcolor{gray}{CONFIDENTIAL}}
\fancyfoot[R]{\small\textcolor{gray}{Page \thepage\ of \pageref{LastPage}}}
\renewcommand{\headrulewidth}{0.4pt}
\renewcommand{\footrulewidth}{0.4pt}

% ─── Colours ────────────────────────────────────────────────────────────────
\definecolor{mkmblue}{HTML}{1565C0}
\definecolor{codegreen}{rgb}{0,0.6,0}
\definecolor{codegray}{rgb}{0.5,0.5,0.5}
\definecolor{codepurple}{rgb}{0.58,0,0.82}
\definecolor{backcolour}{rgb}{0.95,0.95,0.92}

% ─── Hyperref setup ─────────────────────────────────────────────────────────
\hypersetup{
    colorlinks=true,
    linkcolor=mkmblue,
    citecolor=mkmblue,
    urlcolor=mkmblue
}

% ─── Code listings ──────────────────────────────────────────────────────────
\lstdefinestyle{mystyle}{
    backgroundcolor=\color{backcolour},
    commentstyle=\color{codegreen},
    keywordstyle=\color{magenta},
    numberstyle=\tiny\color{codegray},
    stringstyle=\color{codepurple},
    basicstyle=\ttfamily\footnotesize,
    breakatwhitespace=false,
    breaklines=true,
    captionpos=b,
    keepspaces=true,
    numbers=left,
    numbersep=5pt,
    showspaces=false,
    showstringspaces=false,
    showtabs=false,
    tabsize=2
}
\lstset{style=mystyle}

% ─── Section formatting ────────────────────────────────────────────────────
\setcounter{tocdepth}{3}
\setcounter{secnumdepth}{3}

% ─── Shared macros ──────────────────────────────────────────────────────────
\newcommand{\mkmtitlepage}{%
    \begin{titlepage}
    \begin{center}
    \vspace*{1cm}
    \begin{figure}[H]\centering
        \includegraphics[width=0.4\textwidth]{../shared/Logos/MKM_logo.png}
    \end{figure}
    \vspace{1.5cm}
    {\Huge\bfseries\textcolor{mkmblue}{\doctitle}\par}
    \vspace{0.5cm}
    {\Large\textcolor{gray}{\docsubtitle}\par}
    \vspace{2cm}
    {\large MKM Research Labs\par}
    \vspace{0.5cm}
    {\large Version \docversion{} --- \docdate\par}
    \vspace{0.5cm}
    {\large \docauthor\par}
    \vspace{2cm}
    \textcolor{red}{\textbf{CONFIDENTIAL --- PROPRIETARY}}
    \end{center}
    \end{titlepage}
}

\newcommand{\mkmlegalpage}{%
    \clearpage
    \thispagestyle{empty}
    \begin{center}
    \large\textbf{Legal Notice}
    \vspace{1em}

    \textbf{PROPRIETARY AND CONFIDENTIAL}
    \end{center}

    \noindent This document, including all algorithms, models, methodology, formulas,
    calculations, and intellectual concepts contained herein, constitutes the exclusive
    intellectual property and confidential trade secrets of MKM Research Labs (``MKM'').

    \vspace{0.5em}
    \noindent Any usage, reproduction, distribution, modification, or disclosure of this
    document or any portion thereof, without the express prior written authorisation from
    MKM Research Labs is strictly prohibited and constitutes an infringement of intellectual
    property rights.

    \vspace{0.5em}
    \noindent The document is provided on an ``as is'' basis. MKM Research Labs makes no
    representations or warranties of any kind, express or implied, regarding its accuracy,
    reliability, or suitability for any particular purpose.

    \vspace{1em}
    \begin{center}
    All rights reserved. \copyright\ 2021--2026 MKM Research Labs.

    \vspace{0.5em}
    Governed by the laws of the United Kingdom.
    \end{center}
    \clearpage
}

\newcommand{\mkmhistorypage}[1]{%
    \section{Document History}
    \begin{table}[ht]
        \centering
        \begin{tabular}{l|l|l|l}
            \textbf{Date} & \textbf{Version} & \textbf{Description} & \textbf{Author} \\
            \hline
            #1
        \end{tabular}
        \label{tab:revision_history}
    \end{table}
    \clearpage
}

% ─── Bank validation protocol section commands ──────────────────────────────
% Use these to ensure every model doc has the required sections.
% They are just styled \section calls — content is provided by each doc.

\newcommand{\mkmexecsummary}{\section{Executive Summary}}
\newcommand{\mkmpurpose}{\section{Model Purpose and Scope}}
\newcommand{\mkmframework}{\section{Mathematical Framework}}
\newcommand{\mkminputs}{\section{Input Parameters and Calibration}}
\newcommand{\mkmimplementation}{\section{Implementation Details}}
\newcommand{\mkmsensitivity}{\section{Sensitivity Analysis}}
\newcommand{\mkmvalidation}{\section{Validation and Backtesting}}
\newcommand{\mkmtestresults}{\subsection{Automated Test Results}\input{test_results}}
\newcommand{\mkmlimitations}{\section{Model Limitations and Known Weaknesses}}

\newcommand{\mkmgovernance}{\section{Governance Validation Questions}
\noindent The following nine questions are mandated by Chapter~5 of the
\emph{Model Risk Governance for Vendors} handbook and must be explicitly
addressed for every model in the inventory. Response status is tracked in the
Model Governance Dashboard (MRC portal). \par\medskip}

\newcommand{\mkmvqitem}[3]{\subsection{Q#1 --- #2}
\noindent\textit{``#3''}\par\medskip\noindent}

\newcommand{\mkmchangelog}{\section{Change History}}
  (before \begin{document})
%
% Required commands before \input:
%   \newcommand{\doctitle}{...}
%   \newcommand{\docsubtitle}{...}
%   \newcommand{\docversion}{1.0}
%   \newcommand{\docdate}{13-February-2026}
%   \newcommand{\docauthor}{David K Kelly}

% ─── Packages ───────────────────────────────────────────────────────────────
\usepackage[utf8]{inputenc}
\usepackage{amsmath,amsfonts,amssymb}
\usepackage{listings}
\usepackage{color,xcolor}
\usepackage{graphicx}
\usepackage{hyperref}
\usepackage[toc,page]{appendix}
\usepackage{tabto}
\usepackage{float}
\usepackage{comment}
\usepackage[en-GB]{datetime2}
\usepackage{setspace}
\usepackage{subcaption}
\usepackage{booktabs}
\usepackage{longtable}
\usepackage{adjustbox}
\usepackage{enumitem}
\usepackage{fancyhdr}
\usepackage{lastpage}

% ─── Page geometry ──────────────────────────────────────────────────────────
\onehalfspacing
\usepackage{geometry}
\geometry{a4paper, total={170mm,247mm}, left=20mm, top=25mm, bottom=25mm}

% ─── Headers and footers ───────────────────────────────────────────────────
\pagestyle{fancy}
\fancyhf{}
\fancyhead[L]{\small\textcolor{gray}{MKM Research Labs}}
\fancyhead[R]{\small\textcolor{gray}{\doctitle{} --- v\docversion}}
\fancyfoot[C]{\small\textcolor{gray}{CONFIDENTIAL}}
\fancyfoot[R]{\small\textcolor{gray}{Page \thepage\ of \pageref{LastPage}}}
\renewcommand{\headrulewidth}{0.4pt}
\renewcommand{\footrulewidth}{0.4pt}

% ─── Colours ────────────────────────────────────────────────────────────────
\definecolor{mkmblue}{HTML}{1565C0}
\definecolor{codegreen}{rgb}{0,0.6,0}
\definecolor{codegray}{rgb}{0.5,0.5,0.5}
\definecolor{codepurple}{rgb}{0.58,0,0.82}
\definecolor{backcolour}{rgb}{0.95,0.95,0.92}

% ─── Hyperref setup ─────────────────────────────────────────────────────────
\hypersetup{
    colorlinks=true,
    linkcolor=mkmblue,
    citecolor=mkmblue,
    urlcolor=mkmblue
}

% ─── Code listings ──────────────────────────────────────────────────────────
\lstdefinestyle{mystyle}{
    backgroundcolor=\color{backcolour},
    commentstyle=\color{codegreen},
    keywordstyle=\color{magenta},
    numberstyle=\tiny\color{codegray},
    stringstyle=\color{codepurple},
    basicstyle=\ttfamily\footnotesize,
    breakatwhitespace=false,
    breaklines=true,
    captionpos=b,
    keepspaces=true,
    numbers=left,
    numbersep=5pt,
    showspaces=false,
    showstringspaces=false,
    showtabs=false,
    tabsize=2
}
\lstset{style=mystyle}

% ─── Section formatting ────────────────────────────────────────────────────
\setcounter{tocdepth}{3}
\setcounter{secnumdepth}{3}

% ─── Shared macros ──────────────────────────────────────────────────────────
\newcommand{\mkmtitlepage}{%
    \begin{titlepage}
    \begin{center}
    \vspace*{1cm}
    \begin{figure}[H]\centering
        \includegraphics[width=0.4\textwidth]{../shared/Logos/MKM_logo.png}
    \end{figure}
    \vspace{1.5cm}
    {\Huge\bfseries\textcolor{mkmblue}{\doctitle}\par}
    \vspace{0.5cm}
    {\Large\textcolor{gray}{\docsubtitle}\par}
    \vspace{2cm}
    {\large MKM Research Labs\par}
    \vspace{0.5cm}
    {\large Version \docversion{} --- \docdate\par}
    \vspace{0.5cm}
    {\large \docauthor\par}
    \vspace{2cm}
    \textcolor{red}{\textbf{CONFIDENTIAL --- PROPRIETARY}}
    \end{center}
    \end{titlepage}
}

\newcommand{\mkmlegalpage}{%
    \clearpage
    \thispagestyle{empty}
    \begin{center}
    \large\textbf{Legal Notice}
    \vspace{1em}

    \textbf{PROPRIETARY AND CONFIDENTIAL}
    \end{center}

    \noindent This document, including all algorithms, models, methodology, formulas,
    calculations, and intellectual concepts contained herein, constitutes the exclusive
    intellectual property and confidential trade secrets of MKM Research Labs (``MKM'').

    \vspace{0.5em}
    \noindent Any usage, reproduction, distribution, modification, or disclosure of this
    document or any portion thereof, without the express prior written authorisation from
    MKM Research Labs is strictly prohibited and constitutes an infringement of intellectual
    property rights.

    \vspace{0.5em}
    \noindent The document is provided on an ``as is'' basis. MKM Research Labs makes no
    representations or warranties of any kind, express or implied, regarding its accuracy,
    reliability, or suitability for any particular purpose.

    \vspace{1em}
    \begin{center}
    All rights reserved. \copyright\ 2021--2026 MKM Research Labs.

    \vspace{0.5em}
    Governed by the laws of the United Kingdom.
    \end{center}
    \clearpage
}

\newcommand{\mkmhistorypage}[1]{%
    \section{Document History}
    \begin{table}[ht]
        \centering
        \begin{tabular}{l|l|l|l}
            \textbf{Date} & \textbf{Version} & \textbf{Description} & \textbf{Author} \\
            \hline
            #1
        \end{tabular}
        \label{tab:revision_history}
    \end{table}
    \clearpage
}

% ─── Bank validation protocol section commands ──────────────────────────────
% Use these to ensure every model doc has the required sections.
% They are just styled \section calls — content is provided by each doc.

\newcommand{\mkmexecsummary}{\section{Executive Summary}}
\newcommand{\mkmpurpose}{\section{Model Purpose and Scope}}
\newcommand{\mkmframework}{\section{Mathematical Framework}}
\newcommand{\mkminputs}{\section{Input Parameters and Calibration}}
\newcommand{\mkmimplementation}{\section{Implementation Details}}
\newcommand{\mkmsensitivity}{\section{Sensitivity Analysis}}
\newcommand{\mkmvalidation}{\section{Validation and Backtesting}}
\newcommand{\mkmtestresults}{\subsection{Automated Test Results}\input{test_results}}
\newcommand{\mkmlimitations}{\section{Model Limitations and Known Weaknesses}}

\newcommand{\mkmgovernance}{\section{Governance Validation Questions}
\noindent The following nine questions are mandated by Chapter~5 of the
\emph{Model Risk Governance for Vendors} handbook and must be explicitly
addressed for every model in the inventory. Response status is tracked in the
Model Governance Dashboard (MRC portal). \par\medskip}

\newcommand{\mkmvqitem}[3]{\subsection{Q#1 --- #2}
\noindent\textit{``#3''}\par\medskip\noindent}

\newcommand{\mkmchangelog}{\section{Change History}}
  (before \begin{document})
%
% Required commands before \input:
%   \newcommand{\doctitle}{...}
%   \newcommand{\docsubtitle}{...}
%   \newcommand{\docversion}{1.0}
%   \newcommand{\docdate}{13-February-2026}
%   \newcommand{\docauthor}{David K Kelly}

% ─── Packages ───────────────────────────────────────────────────────────────
\usepackage[utf8]{inputenc}
\usepackage{amsmath,amsfonts,amssymb}
\usepackage{listings}
\usepackage{color,xcolor}
\usepackage{graphicx}
\usepackage{hyperref}
\usepackage[toc,page]{appendix}
\usepackage{tabto}
\usepackage{float}
\usepackage{comment}
\usepackage[en-GB]{datetime2}
\usepackage{setspace}
\usepackage{subcaption}
\usepackage{booktabs}
\usepackage{longtable}
\usepackage{adjustbox}
\usepackage{enumitem}
\usepackage{fancyhdr}
\usepackage{lastpage}

% ─── Page geometry ──────────────────────────────────────────────────────────
\onehalfspacing
\usepackage{geometry}
\geometry{a4paper, total={170mm,247mm}, left=20mm, top=25mm, bottom=25mm}

% ─── Headers and footers ───────────────────────────────────────────────────
\pagestyle{fancy}
\fancyhf{}
\fancyhead[L]{\small\textcolor{gray}{MKM Research Labs}}
\fancyhead[R]{\small\textcolor{gray}{\doctitle{} --- v\docversion}}
\fancyfoot[C]{\small\textcolor{gray}{CONFIDENTIAL}}
\fancyfoot[R]{\small\textcolor{gray}{Page \thepage\ of \pageref{LastPage}}}
\renewcommand{\headrulewidth}{0.4pt}
\renewcommand{\footrulewidth}{0.4pt}

% ─── Colours ────────────────────────────────────────────────────────────────
\definecolor{mkmblue}{HTML}{1565C0}
\definecolor{codegreen}{rgb}{0,0.6,0}
\definecolor{codegray}{rgb}{0.5,0.5,0.5}
\definecolor{codepurple}{rgb}{0.58,0,0.82}
\definecolor{backcolour}{rgb}{0.95,0.95,0.92}

% ─── Hyperref setup ─────────────────────────────────────────────────────────
\hypersetup{
    colorlinks=true,
    linkcolor=mkmblue,
    citecolor=mkmblue,
    urlcolor=mkmblue
}

% ─── Code listings ──────────────────────────────────────────────────────────
\lstdefinestyle{mystyle}{
    backgroundcolor=\color{backcolour},
    commentstyle=\color{codegreen},
    keywordstyle=\color{magenta},
    numberstyle=\tiny\color{codegray},
    stringstyle=\color{codepurple},
    basicstyle=\ttfamily\footnotesize,
    breakatwhitespace=false,
    breaklines=true,
    captionpos=b,
    keepspaces=true,
    numbers=left,
    numbersep=5pt,
    showspaces=false,
    showstringspaces=false,
    showtabs=false,
    tabsize=2
}
\lstset{style=mystyle}

% ─── Section formatting ────────────────────────────────────────────────────
\setcounter{tocdepth}{3}
\setcounter{secnumdepth}{3}

% ─── Shared macros ──────────────────────────────────────────────────────────
\newcommand{\mkmtitlepage}{%
    \begin{titlepage}
    \begin{center}
    \vspace*{1cm}
    \begin{figure}[H]\centering
        \includegraphics[width=0.4\textwidth]{../shared/Logos/MKM_logo.png}
    \end{figure}
    \vspace{1.5cm}
    {\Huge\bfseries\textcolor{mkmblue}{\doctitle}\par}
    \vspace{0.5cm}
    {\Large\textcolor{gray}{\docsubtitle}\par}
    \vspace{2cm}
    {\large MKM Research Labs\par}
    \vspace{0.5cm}
    {\large Version \docversion{} --- \docdate\par}
    \vspace{0.5cm}
    {\large \docauthor\par}
    \vspace{2cm}
    \textcolor{red}{\textbf{CONFIDENTIAL --- PROPRIETARY}}
    \end{center}
    \end{titlepage}
}

\newcommand{\mkmlegalpage}{%
    \clearpage
    \thispagestyle{empty}
    \begin{center}
    \large\textbf{Legal Notice}
    \vspace{1em}

    \textbf{PROPRIETARY AND CONFIDENTIAL}
    \end{center}

    \noindent This document, including all algorithms, models, methodology, formulas,
    calculations, and intellectual concepts contained herein, constitutes the exclusive
    intellectual property and confidential trade secrets of MKM Research Labs (``MKM'').

    \vspace{0.5em}
    \noindent Any usage, reproduction, distribution, modification, or disclosure of this
    document or any portion thereof, without the express prior written authorisation from
    MKM Research Labs is strictly prohibited and constitutes an infringement of intellectual
    property rights.

    \vspace{0.5em}
    \noindent The document is provided on an ``as is'' basis. MKM Research Labs makes no
    representations or warranties of any kind, express or implied, regarding its accuracy,
    reliability, or suitability for any particular purpose.

    \vspace{1em}
    \begin{center}
    All rights reserved. \copyright\ 2021--2026 MKM Research Labs.

    \vspace{0.5em}
    Governed by the laws of the United Kingdom.
    \end{center}
    \clearpage
}

\newcommand{\mkmhistorypage}[1]{%
    \section{Document History}
    \begin{table}[ht]
        \centering
        \begin{tabular}{l|l|l|l}
            \textbf{Date} & \textbf{Version} & \textbf{Description} & \textbf{Author} \\
            \hline
            #1
        \end{tabular}
        \label{tab:revision_history}
    \end{table}
    \clearpage
}

% ─── Bank validation protocol section commands ──────────────────────────────
% Use these to ensure every model doc has the required sections.
% They are just styled \section calls — content is provided by each doc.

\newcommand{\mkmexecsummary}{\section{Executive Summary}}
\newcommand{\mkmpurpose}{\section{Model Purpose and Scope}}
\newcommand{\mkmframework}{\section{Mathematical Framework}}
\newcommand{\mkminputs}{\section{Input Parameters and Calibration}}
\newcommand{\mkmimplementation}{\section{Implementation Details}}
\newcommand{\mkmsensitivity}{\section{Sensitivity Analysis}}
\newcommand{\mkmvalidation}{\section{Validation and Backtesting}}
\newcommand{\mkmtestresults}{\subsection{Automated Test Results}\input{test_results}}
\newcommand{\mkmlimitations}{\section{Model Limitations and Known Weaknesses}}

\newcommand{\mkmgovernance}{\section{Governance Validation Questions}
\noindent The following nine questions are mandated by Chapter~5 of the
\emph{Model Risk Governance for Vendors} handbook and must be explicitly
addressed for every model in the inventory. Response status is tracked in the
Model Governance Dashboard (MRC portal). \par\medskip}

\newcommand{\mkmvqitem}[3]{\subsection{Q#1 --- #2}
\noindent\textit{``#3''}\par\medskip\noindent}

\newcommand{\mkmchangelog}{\section{Change History}}

% T1 so the underscores in code identifiers print as underscores.
\usepackage[T1]{fontenc}
\usepackage{lmodern}
% Generated by tools.wolfram_oracle -- do not edit.
\newcommand{\bmGenerated}{2026-10-27}
\newcommand{\bmCodeCommit}{\texttt{d8acf2c90152}}
\newcommand{\bmCodeDirty}{}
\newcommand{\bmNumerics}{numpy 1.26.4, scipy 1.15.1, mpmath 1.4.1}
\newcommand{\bmWolfram}{15.0.0 for Mac OS X ARM (64-bit) (May 26, 2026)}
\newcommand{\bmWolframVia}{ via live kernel session (wolframclient)}
\newcommand{\bmWolframConfirms}{Wolfram and the 50-digit reference agree to $4.9\times10^{-30}$, so the reference is confirmed by a second, independent implementation.}
\newcommand{\bmCases}{161}
\newcommand{\bmChecks}{13}
\newcommand{\bmOutcome}{all pass}
\newcommand{\bmTermLow}{0.10}
\newcommand{\bmTermHigh}{17.3}
\newcommand{\bmPrecision}{$1.2\times10^{-16}$}
\newcommand{\bmPrecisionBelow}{0.0001}
\newcommand{\bmStructure}{$2.6\times10^{-16}$}
\newcommand{\bmNbGap}{$1.8\times10^{-10}$}
\newcommand{\bmImpliedP}{0.29}

% The shared history table sets every column to its natural width; wrap the description.
\renewcommand{\mkmhistorypage}[1]{%
    \section{Document History}
    \begin{table}[ht]
        \centering
        \begin{tabular}{l|l|p{9cm}|l}
            \textbf{Date} & \textbf{Version} & \textbf{Description} & \textbf{Author} \\
            \hline
            #1
        \end{tabular}
    \end{table}
}


\begin{document}
% Long code paths and identifiers cannot hyphenate; let the spacing stretch instead.
\sloppy

\mkmtitlepage
\mkmlegalpage

\tableofcontents
\clearpage

% =========================================================================
\mkmexecsummary
% =========================================================================

This document benchmarks the Event Frequency Model (MKM-EF-001), which turns a
gauge's event arrival rate $\lambda$ and its per-event flood probability $p$ into
the annual flood probability every PRS is priced from, and into the multi-year term
structure the gauge PRS pricer reads. It asks whether the code computes what its
model document (\texttt{event\_frequency\_plan.pdf}) says, and whether what the
document says holds. Version~0.1 tested the code against v2.3 of the document and
found the term structure understating multi-year risk and annualisation losing
precision at small $\lambda p$. Both are fixed, and the document is issued as v2.4,
which this version tests.

\paragraph{The code computes its documented formulae.} Over \bmCases{} cases ---
arrival rates from 0.5 to 50 a year, per-event probabilities from $10^{-6}$ to 0.5,
tenors of 1 to 10 years with a constant and a growing rate, three loss catalogues
and seven count series --- every quantity agrees with a 50-digit implementation of
the document's equations to better than $10^{-12}$ (Table~\ref{tab:bm-checks}).
\bmWolframConfirms{} Wolfram's own maximisation of the Negative Binomial likelihood
finds nothing higher than the model's fit.

\paragraph{The term structure is exact.} Each year now contributes its exceedance
intensity $\lambda_y p$, so the probability of a flood within a tenor is
$1-\exp(-p\sum_y\lambda_y)$ at every point, to \bmStructure{}
(Table~\ref{tab:bm-term}). v2.3's code compounded each year's annual
\emph{probability} as though it were an intensity, which understated that
probability by \bmTermLow\% to \bmTermHigh\%, worst at one year and high
$\lambda p$. The correction raises every multi-year probability, and every PRS
spread priced from the term structure.

\paragraph{Annualisation keeps its precision.} Evaluated as $-\mathrm{expm1}(-\lambda p)$,
the annual probability is exact to \bmPrecision{} below $\lambda p=\bmPrecisionBelow$,
where $1-\exp(-x)$ lost $10^{-10}$.

\paragraph{The document's claims hold} (Table~\ref{tab:bm-claims}): the closed form
of \S5.9 is an identity (Wolfram proves it symbolically), the gate statistic's mean
$|z|$ is $\sqrt{2/\pi}=0.798$, the sampling errors of \S7.2 scale as $1/\sqrt n$, and
the Negative Binomial search reaches its maximum.

\paragraph{Evidence.} \bmChecks{} checks, outcome: \bmOutcome{}
(Table~\ref{tab:bm-checks}).
\begin{itemize}[leftmargin=*,nosep]
\item Wolfram: \bmWolfram{}\bmWolframVia.
\item Code under test, at commit \bmCodeCommit{}\bmCodeDirty.
\item Numerics: \bmNumerics.
\item Generated \bmGenerated{} by \texttt{python -m tools.wolfram\_oracle ef001}.
\end{itemize}

% =========================================================================
\mkmpurpose
% =========================================================================

\subsection{What is benchmarked}
In \texttt{src/models/frequency}: \texttt{annual\_exceedance\_probability},
\texttt{annual\_exceedance\_rate} and \texttt{annual\_hazard\_by\_year}
(\texttt{annualise.py}); \texttt{term\_exceedance\_probability} and the rate
processes (\texttt{rate\_process.py}); \texttt{analytic\_average\_annual\_loss} and
\texttt{loss\_standard\_error} (\texttt{ylt/\_reconcile.py}); and the count
families, \texttt{fit\_poisson}, \texttt{fit\_negbin} and \texttt{dispersion\_test}
(\texttt{families/}). With them, \texttt{compute\_term\_structure} in
\texttt{src/models/hazard/gev.py}, which the per-year hazards feed. Each file is
loaded on its own and called as production calls it.

\subsection{What is not}
The Monte Carlo year simulation is checked by its own reconciliation gate, whose
statistic is tested here; the simulation's draws are not. The calibration of
$\lambda$ from gauge records, which the document's \S6 and \S8 identify as the
model's open problem, is not a calculation an oracle can check. This is a
benchmark of the implementation (validation question Q8), not independent
validation (Q9).

% =========================================================================
\mkmframework
% =========================================================================

With $N\sim\text{Poisson}(\lambda)$ events a year, each flooding independently with
probability $p$, the document's \S5.9 gives
\begin{equation}
\mathbb{E}\bigl[(1-p)^N\bigr]=e^{-\lambda p},\qquad
P(\text{at least one flood in a year})=1-e^{-\lambda p},
\end{equation}
and the aggregate rate $\lambda p$, which is never smaller. Over a tenor of $T$
years with rates $\lambda_0,\dots,\lambda_{T-1}$ (constant, or growing
geometrically as $\lambda_0(1+g)^y$), the exact tenor probability is
$1-\exp\bigl(-p\sum_y\lambda_y\bigr)$.

\paragraph{The term structure} (v2.4, \S5.10). The hazard builder hands
\texttt{compute\_term\_structure} each year's intensity $\lambda_y p$, which it
accumulates, so each point is the exact tenor probability. v2.3's code handed it the
year's \emph{probability} $q_y=1-e^{-\lambda_y p}$ instead, giving
$1-\exp\bigl(-\sum_y q_y\bigr)$; since $q_y<\lambda_y p$ that always understated, by
more as $\lambda p$ grows, and most relative to the answer at one year.

\paragraph{Compound loss.} With event weights $w_k$ and losses $L_k$, the annual
loss is compound Poisson, with mean $\lambda\sum_k w_kL_k$ and, over $n$ simulated
years, standard error $\sqrt{\lambda\sum_k w_kL_k^2/n}$.

\paragraph{Count families.} The dispersion statistic $(n-1)s^2/\bar x$ is
$\chi^2_{n-1}$ under Poisson; the Negative Binomial is fitted with its mean at
$\bar x$ and its dispersion $\alpha=1/r$ by bounded search, which is the joint
maximum likelihood because the Negative Binomial's estimating equation for the mean
is solved by $\bar x$ at every $r$.

% =========================================================================
\section{Benchmark Method}
% =========================================================================

The equations above are evaluated with \texttt{mpmath} at 50 significant digits, from
the document rather than the code. Wolfram evaluates them independently from each
input's exact value at 60-digit precision: the compound loss through
\texttt{CompoundPoissonDistribution} of an \texttt{EmpiricalDistribution} (its rate
scaled by the weights' exact total, which \texttt{EmpiricalDistribution} would
otherwise renormalise away), the $p$-values through \texttt{ChiSquareDistribution},
and the Negative Binomial maximum by its own \texttt{FindMaximum} over both
parameters, with the Poisson limit as the supremum for under-dispersed data. It
proves the identity of \S5.9 with \texttt{Expectation} and \texttt{Simplify}.
Every annual probability is judged against the $10^{-12}$ tolerance, down to
$\lambda p=5\times10^{-7}$.

% =========================================================================
\mkmvalidation
% =========================================================================

% Generated by tools.wolfram_oracle -- do not edit.
\begin{table}[H]\centering\small
\caption{Benchmark checks}\label{tab:bm-checks}
\begin{adjustbox}{max width=\textwidth}
\begin{tabular}{p{88mm}p{22mm}rl}
\toprule
Check & Result & Tolerance & Outcome \\
\midrule
Annual probability $1-e^{-\lambda p}$, $\lambda p\ge0.0001$: model vs reference & $5.7\times10^{-14}$ & $1.0\times10^{-12}$ & pass \\
Aggregate rate $\lambda p$: model vs reference & $5.6\times10^{-17}$ & $1.0\times10^{-12}$ & pass \\
Tenor probability $1-e^{-p\sum\lambda_y}$: model vs reference & $6.1\times10^{-15}$ & $1.0\times10^{-12}$ & pass \\
Compound-Poisson average annual loss: model vs reference & $2.3\times10^{-16}$ & $1.0\times10^{-12}$ & pass \\
Its standard error over $n$ years: model vs reference & $4.6\times10^{-17}$ & $1.0\times10^{-12}$ & pass \\
Dispersion statistic $(n-1)s^2/\bar x$: model vs reference & $1.4\times10^{-16}$ & $1.0\times10^{-12}$ & pass \\
Over-dispersion $p$-value ($\chi^2$ survival): model vs reference & $9.4\times10^{-15}$ & $1.0\times10^{-10}$ & pass \\
Under-dispersion $p$-value ($\chi^2$ CDF): model vs reference & $5.5\times10^{-16}$ & $1.0\times10^{-10}$ & pass \\
Poisson log-likelihood: model vs reference & $1.3\times10^{-15}$ & $1.0\times10^{-12}$ & pass \\
Negative Binomial log-likelihood at its maximum: model vs reference & $1.3\times10^{-12}$ & $1.0\times10^{-8}$ & pass \\
Wolfram agrees with the reference (closed forms, $\chi^2$, Poisson) & $4.9\times10^{-30}$ & $1.0\times10^{-12}$ & pass \\
Wolfram's own Negative Binomial maximum finds no higher likelihood & 0 & $1.0\times10^{-8}$ & pass \\
\bottomrule
\end{tabular}
\end{adjustbox}
\end{table}


\subsection{The model document's claims}
% Generated by tools.wolfram_oracle -- do not edit.
\begin{table}[H]\centering\small
\caption{The model document's claims, and the code's own conventions, tested}\label{tab:bm-claims}
\begin{adjustbox}{max width=\textwidth}
\begin{tabular}{p{36mm}p{44mm}p{56mm}l}
\toprule
Claim & Stated & Benchmark finds & Outcome \\
\midrule
\S5.9 closed form & $E[(1-p)^N]=e^{-\lambda p}$; aggregate $\ge$ occurrence & identity confirmed symbolically; aggregate $\ge$ occurrence in every case & met \\
\S5.9 gate statistic & mean $|z|=\sqrt{2/\pi}=0.798$ & $0.797885$ & met \\
\S7.2 sampling error & 5, 1.6, 0.5, 0.16\% at $10^3$..$10^6$ years & $1/\sqrt n$ from the first gives 5.00\%, 1.58\%, 0.50\%, 0.16\% & met \\
Term structure (annualise, gev) & annual probability used as the per-year hazard & tenor probability understated by 0.10--17.3\% & finding \\
Annualisation precision & $1-\exp(-x)$ & $1.0\times10^{-10}$ relative below $\lambda p=0.0001$; $-\mathrm{expm1}(-x)$ is exact & finding \\
Negative Binomial search & bounded scalar search on $\alpha$ & log-likelihood within $1.8\times10^{-10}$ of the maximum & met \\
\bottomrule
\end{tabular}
\end{adjustbox}
\end{table}


\noindent The \S7.2 errors imply a relative error $\sqrt{(1-q)/(nq)}$ with
$q\approx\bmImpliedP$: the table quotes one subject's figures, not a bound.

\subsection{The multi-year term structure}
% Generated by tools.wolfram_oracle -- do not edit.
\begin{table}[H]\centering\small
\caption{The multi-year exceedance probability: the term structure now, and v2.3's convention}\label{tab:bm-term}
\begin{adjustbox}{max width=\textwidth}
\begin{tabular}{lrrrrrrr}
\toprule
Rate & $\lambda_0$ & $p$ & $T$ & Exact & Now & v2.3 & v2.3 gap \\
\midrule
const & 2 & 0.001 & 1 & 0.001998 & 0.001998 & 0.001996 & -0.10\% \\
const & 2 & 0.001 & 10 & 0.019801 & 0.019801 & 0.019782 & -0.10\% \\
const & 2 & 0.01 & 1 & 0.019801 & 0.019801 & 0.019607 & -0.98\% \\
const & 2 & 0.01 & 10 & 0.181269 & 0.181269 & 0.179641 & -0.90\% \\
const & 2 & 0.05 & 1 & 0.095163 & 0.095163 & 0.090775 & -4.61\% \\
const & 2 & 0.05 & 10 & 0.632121 & 0.632121 & 0.613887 & -2.88\% \\
const & 10 & 0.001 & 1 & 0.009950 & 0.009950 & 0.009901 & -0.50\% \\
const & 10 & 0.001 & 10 & 0.095163 & 0.095163 & 0.094712 & -0.47\% \\
const & 10 & 0.01 & 1 & 0.095163 & 0.095163 & 0.090775 & -4.61\% \\
const & 10 & 0.01 & 10 & 0.632121 & 0.632121 & 0.613887 & -2.88\% \\
const & 10 & 0.05 & 1 & 0.393469 & 0.393469 & 0.325288 & -17.33\% \\
const & 10 & 0.05 & 10 & 0.993262 & 0.993262 & 0.980448 & -1.29\% \\
trend & 5 & 0.001 & 1 & 0.004988 & 0.004988 & 0.004975 & -0.25\% \\
trend & 5 & 0.001 & 10 & 0.060953 & 0.060953 & 0.060764 & -0.31\% \\
trend & 5 & 0.01 & 1 & 0.048771 & 0.048771 & 0.047600 & -2.40\% \\
trend & 5 & 0.01 & 10 & 0.466819 & 0.466819 & 0.456195 & -2.28\% \\
trend & 5 & 0.05 & 1 & 0.221199 & 0.221199 & 0.198443 & -10.29\% \\
trend & 5 & 0.05 & 10 & 0.956910 & 0.956910 & 0.932189 & -2.58\% \\
\bottomrule
\end{tabular}
\end{adjustbox}
\end{table}


% =========================================================================
\mkmlimitations
% =========================================================================

\begin{enumerate}[leftmargin=*]
\item \textbf{Benchmark, not validation.} The code is shown to compute the
  documented formulae; whether $\lambda$ is right for a gauge is the question the
  document's own \S8 raises and leaves open.
\item \textbf{Stored term structures.} Term structures already written by the hazard
  builder, and any spreads priced from them, carry v2.3's understatement until the
  hazard curves are rebuilt.
\item \textbf{Naming.} The hazard curve still stores the annual \emph{probability}
  under \texttt{annual\_hazard\_rate\_*}; PRS pricing reads it as a probability, so
  the value is right and only the name misleads.
\end{enumerate}

\mkmhistorypage{26-Oct-2026 & 0.1 & First issue: tests v2.3; finds the term structure understated & David K Kelly \\
27-Oct-2026 & 0.2 & After the fixes: tests v2.4 & David K Kelly \\}

\end{document}
